\subsection{An invariant relative category with a non-surjective comparison}

\begin{claim}{C.17}{Proved}{The crown: the exact ind-object and every test category}
Let $M$ be the crown poset on $\mathbb Z/4$, with nonidentity
arrows $0\to1,0\to3,2\to1,2\to3$, and mark every arrow.
For all small categories $X$, the specified path comparison at
$a=b=0$ identifies
\[
 H_{0,0}(X)\cong
 \{f:\pi_0X\to\mathbb Z\mid f\text{ has finite image}\},
 \qquad
 \operatorname{Map}_{L(M,M)}(0,0)\simeq\mathbb Z_{\mathrm{disc}}.
                                                               \tag{C.17.1}
\]
Under these identifications $c_X$ is inclusion into the set of
all functions $\pi_0X\to\mathbb Z$. In particular $H_{0,0}$
is globally $W_\infty$-invariant, but $c_X$ is not surjective
for the countably infinite discrete category. The same invariance
holds at every pair of endpoints. The actual symmetric path
colimit has no categorical homotopy equivalence to a local
category, despite this invariance of its ind-object. Its actual
padding maps $q_L,q_R$ are weak equivalences which have no
inverse in $\mathrm{Cat}_{\mathrm{hot}}$.
\end{claim}
\begin{proof}
Use the block $(+,-)$ and let $K_n$ be the symmetric stage,
indexed by the fence vertices $-2n,\ldots,2n$, both endpoints
zero. Its transitions extend constantly at both ends. Reduction
modulo four from the integer fence, whose arrows leave every even
integer towards its adjacent odd integers, is a categorical
covering of $M$. Every path has a unique lift whose left endpoint
is zero. Each arrow of an endpoint-fixed ladder lifts uniquely
with that same left endpoint: lift the square strip
$I_\tau\times[1]$, whose nerve is contractible, and use the
unique lift of the identity at the left. At the right an identity
can only lift to an identity. Thus the lifted right endpoint is
the same for both rows and has the form $4k$. This defines a
functor
\[
 w_n:K_n\longrightarrow D_n=[-n,n]_{\mathrm{disc}}.
\]
Every generating edge changes its integer value by at most one,
so the bound $|k|\leq n$ follows. The $w_n$ commute strictly
with the actual symmetric transitions.

For $k\in\mathbb Z$, define a lifted geodesic on the full
integer parameter fence by being zero for $i\leq0$, having value
$\operatorname{sgn}(k)i$ for $0\leq i\leq4|k|$, and having
value $4k$ afterwards. Its reduction is an eventually constant
path of winding $k$; the sign reversal retains parity and hence
preserves every oriented generating edge. For $|k|\leq j$ it
belongs to $K_{2j}$. These choices define strictly compatible
functors $g_j:D_j\to K_{2j}$ and $w g=1$ as an ind-morphism.

We prove the other inverse relation with its actual transition
maps. Fix $n$ and let $K_n^k$ be the full subcategory with winding
$k$. There are finitely many such parts, with no arrows between
them. All values of the lifted universal path, and of the chosen
geodesics for these $k$, lie in the finite fence
$J_n=[-4n,4n]$. Each finite fence is categorically contractible:
delete an endpoint by sending it to its unique neighbor; this
retraction is connected to the identity by the incident arrow,
and iteration reaches a point. Postcomposition therefore makes
$\operatorname{Fun}(Y,J_n)$ categorically contractible for
\emph{every} small $Y$. The homotopies have a fixed finite
length depending on $J_n$, with no restriction on $Y$.

Take $Y=K_n^k$. The lifted universal path is a localization word
in $J_n^Y$ from the constant zero diagram to the constant $4k$
diagram. The constant geodesic is another such word. Every arrow
of $J_n^Y$ is marked, and its categorical contraction implies
that its ordinary groupoid completion has exactly one arrow
between any two objects. The two words are therefore equal in
this localization. Apply the proved formula C.20 to
$(J_n,\operatorname{Mor}J_n)$, this $Y$, and the actual
symmetric type sequence. Equality is witnessed at some later
stage by a finite zigzag of natural transformations between the
two functors $Y\to F_{t_N}(J_n;0,4k)$. Projection to $M$
gives precisely a homotopy between the actual transition
$K_n^k\to K_N$ and its factorization through the selected
constant geodesic. Choose one common later stage for the finitely
many $k$. On the coproduct $K_n=\coprod_kK_n^k$, concatenate
these finite zigzags, using identity transformations on the other
parts. The result is a genuine finite categorical homotopy
\[
 K_n\longrightarrow K_N
       \simeq K_n\xrightarrow{w_n}D_n
                     \xrightarrow{g_n}K_{2n}\longrightarrow K_N.
\]
This proves $gw=1$ in $\operatorname{Ind}(\mathrm{Cat}_{\mathrm{hot}})$.
No assertion about a contraction inside a fixed path stage was
needed. The verified equalities at later stages give the inverse
ind-morphisms.

Apply $\operatorname{colim}\pi_0\operatorname{Fun}(X,-)$ to
these inverse morphisms. A functor into $D_n$ is exactly a
function $\pi_0X\to[-n,n]$, and this functor category is
discrete. Their union is the finite-image functions in (C.17.1).
C.20 identifies it with $H_{0,0}(X)$, including infinite and
empty $X$. C.7 identifies the localization mapping space with
$\Omega S^1$, since the crown nerve is the four-cycle. The
lifting endpoint $4k$ is the usual degree of a loop and identifies
its components with $\mathbb Z$. The stagewise classifying map
therefore sends a functor to the same winding function. As the
target mapping space is homotopy discrete, this proves that
$c_X$ is exactly the claimed inclusion. For $X=\mathbb N$
discrete, the function $i\mapsto i$ is absent from its image.
Weak equivalences induce bijections on component sets, proving
the global invariance.

For arbitrary endpoints choose integer lifts $\alpha,\beta$.
The right lift now has the form $\beta+4k$ and the effect set is
the corresponding $\mathbb Z$-torsor. Each finite stage has
finitely many such endpoints. The same finite-fence contraction,
using constants $\alpha$ and $\beta+4k$, proves its ind-object
is the finite-subset presentation of this torsor. This proves
invariance at all endpoint pairs, with the choices used only to
name its elements.

Finally let $P=\operatorname{colim}_nK_n$, and let $P_0$ be
its winding-zero part. It has weakly contractible nerve. A
categorical homotopy equivalence of $P$ with a local category
would make $P$ local, since applying $\operatorname{Fun}(X,-)$
retains the inverse zigzags and their classifying maps. The
local comparison for $X=*$ and $X=P_0$ would then contract
$P_0$ in $\mathrm{Cat}_{\mathrm{hot}}$ to a chosen path $c$.
Let a proposed contraction have length $l$. Evaluate its lifted
objectwise zigzags at an integer parameter beyond the support of
$c$. Each of the $l$ vertical arrows changes that integer value
by at most one. However the excursion path with lifted values
$0$ up to zero, $i$ for $0\leq i\leq4N$, $8N-i$ for
$4N\leq i\leq8N$, and zero afterwards, lies in $P_0$ and has
value $4N$ at $4N$. Choose $4N>l$ beyond the support of $c$.
The resulting distance contradiction excludes the contraction.
This also checks the obstruction underlying PATHS C.P5.

The same distance estimate applies directly to the actual padding
map $q_L:L\to P$. Its image consists of paths whose right half,
at integer parameters $i\geq0$, is constant at zero. Suppose a
functor $r:P\to L$ and a finite zigzag of $l$ natural
transformations witnessed $q_Lr\simeq1_P$. Apply that zigzag to
the winding-zero excursion above. Every intervening path has
winding zero, since every arrow preserves the lifted right
endpoint. The lift of its $q_Lr$ image is consequently zero on
the whole right half. Evaluation at $i=4N$ would join this zero
to the excursion's value $4N$ by $l$ steps, each of size at most
one. Choosing $4N>l$ is a contradiction. Thus $q_L$ has no
right inverse even up to categorical homotopy. Reflection of the
integer parameter, $i\mapsto-i$, is a functor of the integer
fence: it preserves parity and sends each arrow from an even
integer to the corresponding arrow to an adjacent odd integer.
It swaps the two half-line systems and intertwines $q_L$ and
$q_R$. Applying the same argument after this reflection excludes
an inverse for $q_R$ as well. C.4 supplies their specified weak
equivalences.

The finite base, covering rule and excursion formulas are
checked in the accompanying path programs; their uniform bounds
prove the statements for every $N$.
\end{proof}
